\section{The explicit binary correction and the collision interpretation}
\label{sec:rational}
We now convert the exact-saddle probability into a more accessible carrier using the same $u$ as in \eqref{eq:ueq}.

\begin{theorem}[Rational correction]\label{thm:rational}
Let
\begin{equation}\label{eq:Ru}
 \mathcal R(u)=\frac{u(u+2)(u^4+8u^3+22u^2+20u+8)}
                     {4(u^2+4u+2)^2}.
\end{equation}
Then
\begin{equation}\label{eq:explicitratio}
 p_N=e^{-u^2/2-u}\left[1+\frac{\mathcal R(u)}N
                         +O\!\left(\frac{(\log N)^8}{N^2}\right)\right].
\end{equation}
Consequently,
\begin{equation}\label{eq:ratiolimits}
 p_Ne^{u^2/2+u}\to1,\qquad
 \frac N{(\log N)^2}\bigl(p_Ne^{u^2/2+u}-1\bigr)\to\frac14,
 \qquad \frac{\log p_N}{(\log N)^2}\to-\frac12.
\end{equation}
The binary count has the relative carrier
\begin{equation}\label{eq:ccarrier}
 c_N=\sqrt{\frac2{P(u)}}
 \exp\!\left(\frac{Nu(u+1)}{u+2}-\frac u2-\frac{u^2}4\right)
 \left(1+O\!\left(\frac{(\log N)^4}{N}\right)\right).
\end{equation}
\end{theorem}
\begin{proof}
Use $x_0=2N/(u+2)$ and the following rational functions:
\begin{equation}\label{eq:algebra-abbr}
 \begin{aligned}
 D&=u+4+2/u,& E&=2u+6-4/u^2,\\
 F&=(u+1)(u+2),& G&=(u+2)(3u+4),\\
 B_0&=-F/2+1/u.&
 \end{aligned}
\end{equation}
Here $B_0$ is unrelated to the integral abbreviation $B=N-1$. We have $-H_N''(x_0)=D/x_0$ and $H_N'''(x_0)=E/x_0^2$.

The differentiated gamma quotient expansion near $x_0$ gives
\begin{equation}\label{eq:fprimefine}
 f_N'(x)=H_N'(x)-\frac1x-\frac{3q}{x^2}-\frac{2q^2}{x^3}
                          +\frac1{2q}+O(L^5/N^2).
\end{equation}
One can obtain this directly by differentiating
$q\log M-\log\Gamma(q+1)+q(q-1)/(2M)$ with its differentiated remainder, and using $\psi(q+1)=\log q+1/(2q)+O(q^{-2})$. The same analytic-disk method as in Section~\ref{sec:binary}, or fixed-order polygamma estimates, bounds the differentiated remainder and avoids differentiating an unspecified error term.
At $q_0=ux_0/2$, \eqref{eq:fprimefine} yields
\begin{equation}\label{eq:rfine}
 f_N'(x_0)=B_0/x_0+O(L^5/N^2),\qquad
 r=x_0+B_0/D+O(L^3/N).
\end{equation}
The last bound follows from one Newton correction and local concavity: the second-order phase term and the variation of the first correction each contribute $O(L^5/N^2)$ to the derivative equation, then division by $A\asymp L^2/N$ gives $O(L^3/N)$.

Put $H_0(x)=-2N(N-x)/x^2$. At $x=x_0$,
\begin{equation}\label{eq:H0values}
 H_0=-u(u+2)/2,\qquad H_0'=F/x_0,\qquad H_0''=-G/x_0^2.
\end{equation}
Expansion of the exact penalty $h_N=H+J+O(L^6/N^2)$ at $x_0$ gives the cancellation
\begin{equation}\label{eq:h0cancel}
 h_N(x_0)=-u(u+2)/2+u/x_0+O(L^6/N^2).
\end{equation}
For clarity, expand $H=H_0(1-1/N)(1+1/x_0)^{-1}$ and
$J=-x_0(N-x_0)(N-1)/(2M(x_0)^2)$. Their terms of order $x_0^{-1}$ sum to $u/x_0$; the omitted terms are bounded by $O(L^6/N^2)$. Transport to $r$ using \eqref{eq:rfine} adds $FB_0/(Dx_0)$ at the same accuracy.

In the Gaussian correction \eqref{eq:binary-first}, replace $A$ by $D/x_0$, $f_N'''$ by $E/x_0^2$, and the penalty derivatives by \eqref{eq:H0values}. All replacement errors are $O(L^8/N^2)$. Thus
\begin{equation}\label{eq:rational-pre}
 p_Ne^{u(u+2)/2}
 =1+\frac1{x_0}\left[u+\frac{FB_0}D+\frac{F^2-G}{2D}
                              +\frac{EF}{2D^2}\right]+O(L^8/N^2).
\end{equation}
Since $B_0=-F/2+1/u$, the two $F^2$ terms cancel. Rational simplification gives
\[
 u+\frac{FB_0}D+\frac{F^2-G}{2D}+\frac{EF}{2D^2}
 =\frac{u(u^4+8u^3+22u^2+20u+8)}{2P(u)^2}.
\]
Finally $1/x_0=(u+2)/(2N)$, proving \eqref{eq:explicitratio}. Before cancellation individual corrections can have order $L^4/N$; their squares explain the safe $O(L^8/N^2)$ remainder. Cancellation of the coefficient does not justify shrinking that remainder. Since $\mathcal R(u)\sim u^2/4$ and $u\sim L$, \eqref{eq:ratiolimits} follows. Multiply \eqref{eq:explicitratio} by \eqref{eq:bcarrier} to obtain \eqref{eq:ccarrier}.
\end{proof}

\subsection{A Lambert W equivalent}
If $w=W_0(\sqrt{N/2})$ and $v=2w$, then $v^2e^v=2N$. Comparing the increasing defining equations gives
\[
 u=v-2/v+O(v^{-2}),\qquad
 u^2/2+u=v^2/2+v-2+O(v^{-1}).
\]
For example, subtract their logarithms to get
$u+2\log u+\log(1+2/u)=v+2\log v$, then use the mean-value theorem and one substitution. Therefore
\begin{equation}\label{eq:Wratio}
 p_N\sim\exp(2-2w^2-2w).
\end{equation}
Its relative error is only $O(1/\log N)$, much weaker than \eqref{eq:explicitratio}. The factor $e^2$ cannot be discarded by replacing $u(u+2)$ by $u^2$ in the implicit equation.

\subsection{Why the probability has this exponent}
After subtracting the mandatory diagonal ones, a fixed-dimension table is a uniform weak composition of $q$ into $M$ cells. For any particular residual cell $X$,
\begin{equation}\label{eq:celltail}
 \Prob(X\ge k\mid m)=\frac{\fall qk}{\fall{M+q-1}k},\qquad 0\le k\le q.
\end{equation}
Here $\fall ak=a(a-1)\cdots(a-k+1)$, with $\fall a0=1$. The formula follows by subtracting $k$ from that cell and taking the ratio of the two composition counts. Let $D_{\rm bad}$ count diagonal cells with positive residual and $C_{\rm bad}$ count off-diagonal cells with residual at least two. Then
\begin{align}
 \E(D_{\rm bad}\mid m)&=\frac{mq}{M+q-1},\label{eq:bad-diag}\\
 \E(C_{\rm bad}\mid m)&=\frac{(M-m)q(q-1)}{(M+q-1)(M+q-2)}.\label{eq:bad-coll}
\end{align}
The first expectation is zero when $q=0$, and the second is zero when $q<2$; these conventions replace any degenerate rational expression. At an integer $m=x_0+O(\log N)$, these means are $u+o(1)$ and $u^2/2+o(1)$, respectively. Binary means both bad counts vanish, so these expectations explain the two parts of the exponent. They are an interpretation, not a proof of a relative zero-event probability. A total-variation or Gaussian approximation with a merely small absolute error could overwhelm this event. The exact product, strengthened tails and quotient expansion above prove the probability directly; no full collision-distribution theorem is claimed.
